\documentclass[a4paper,12pt,twoside]{book}

%%%%%%%%%%%%%%%%%%%%%%%%%%%%%%%%%%%%%%%%%%%%%%%%%%%%%%%%%%%%%%%%%%%%%%%%%%%%%%
%% § Paquetes adicionales                                                   %%
%%%%%%%%%%%%%%%%%%%%%%%%%%%%%%%%%%%%%%%%%%%%%%%%%%%%%%%%%%%%%%%%%%%%%%%%%%%%%%

% Configuración para XeLaTeX
\usepackage{tocloft}
\setlength{\cftbeforechapskip}{2ex}
\setlength{\cftbeforesecskip}{0.5ex}
\setlength{\cftchapnumwidth}{2em}
\setlength{\cftsecnumwidth}{14mm}
\cftsetindents{chapter}{0em}{\cftchapnumwidth}

\usepackage{fontspec}
\usepackage{xunicode}
\usepackage{xltxtra}
\defaultfontfeatures{Ligatures=TeX,Numbers=OldStyle} % inicial
\setromanfont{DejaVu Sans} % inicial
% \setmonofont{DejaVu Sans Mono}[Scale={0.90}] % inicial
\setmonofont{JuliaMono}[Scale=MatchLowercase]

% \setsansfont{Arial}§

% \setmainfont{DejaVu Math TeX Gyre}
% \setmonofont{DejaVu Math TeX Gyre}[Scale={1.1}]

% Notas: La lista de fuentes disponibles se obtiene con fc-list

\usepackage[spanish]{babel}          % Castellanización.
% \usepackage{ucs}
% \usepackage[utf8]{inputenc}        % Acentos de UTF8
% \usepackage[T1]{fontenc}           % Codificación T1
% \usepackage{graphicx}
% \usepackage{fancyvrb}                % Verbatim extendido
% \usepackage{mathpazo}              % Fuentes semejante a palatino
% \usepackage[scaled=.90]{helvet}
% \usepackage{cmtt}
% \renewcommand{\ttdefault}{cmtt}
\usepackage{a4wide}
% \usepackage{amsmath,amssymb,amstext}
\usepackage{amssymb, amsmath, mathtools}

\usepackage{newunicodechar}
\newunicodechar{⟹}{\ensuremath{\mathnormal\Longrightarrow}}
\usepackage[dvipsnames]{xcolor}

\usepackage{minted}
\usemintedstyle{emacs}
\definecolor{bg}{rgb}{0.95,0.95,0.95}
\newminted{lean}{bgcolor=bg, autogobble=true}

\setminted{fontsize=\small,baselinestretch=1}

\linespread{1.05}                  % Distancia entre líneas
\setlength{\parindent}{2em}        % Indentación de comienzo de párrafo
% \deactivatetilden                  % Elima uso de ~ para la eñe
\raggedbottom                      % No ajusta los espacios verticales.

\usepackage[%
  colorlinks=true,
  urlcolor=blue,
  % pdftex,
  pdfauthor={José A. Alonso <jalonso@us.es>},%
  pdftitle={Calculemus con Lean},%
  pdfstartview=FitH,%
  bookmarks=false]{hyperref}

\setcounter{tocdepth}{1}
\setcounter{secnumdepth}{4}

% No indenta los item
\usepackage{enumitem}
\setlist[itemize]{leftmargin=*}

\usepackage{makeidx}
\makeindex

%%%%%%%%%%%%%%%%%%%%%%%%%%%%%%%%%%%%%%%%%%%%%%%%%%%%%%%%%%%%%%%%%%%%%%%%%%%%%%
%% § Cabeceras                                                              %%
%%%%%%%%%%%%%%%%%%%%%%%%%%%%%%%%%%%%%%%%%%%%%%%%%%%%%%%%%%%%%%%%%%%%%%%%%%%%%%

\usepackage{fancyhdr}

\addtolength{\headheight}{\baselineskip}

\pagestyle{fancy}

\cfoot{}

%\fancyhead{}
%\fancyhead[RE]{\mdseries\sffamily \nouppercase{\leftmark}}
%\fancyhead[LO]{\mdseries\sffamily \nouppercase{\rightmark}}
%\fancyhead[LE,RO]{\mdseries\sffamily \thepage}

\fancyhead{}
\fancyhead[RE]{\mdseries \nouppercase{\leftmark}}
\fancyhead[LO]{\mdseries \nouppercase{\rightmark}}
\fancyhead[LE,RO]{\mdseries \thepage}

%%%%%%%%%%%%%%%%%%%%%%%%%%%%%%%%%%%%%%%%%%%%%%%%%%%%%%%%%%%%%%%%%%%%%%%%%%%%%%
%% § Definiciones                                                           %%
%%%%%%%%%%%%%%%%%%%%%%%%%%%%%%%%%%%%%%%%%%%%%%%%%%%%%%%%%%%%%%%%%%%%%%%%%%%%%%

\input definiciones

%%%%%%%%%%%%%%%%%%%%%%%%%%%%%%%%%%%%%%%%%%%%%%%%%%%%%%%%%%%%%%%%%%%%%%%%%%%%%%
%% § Título                                                                 %%
%%%%%%%%%%%%%%%%%%%%%%%%%%%%%%%%%%%%%%%%%%%%%%%%%%%%%%%%%%%%%%%%%%%%%%%%%%%%%%

\title{
  {\LARGE Calculemus \\
    (Vol. 2: Demostraciones con Lean4)}}
\author{
  \href{http://www.cs.us.es/~jalonso}{José A. Alonso Jiménez}}
\date{\vfill \hrule \vspace*{2mm}
  \begin{tabular}{l}
      \href{http://www.cs.us.es/glc}
           {Grupo de Lógica Computacional} \\
      \href{http://www.cs.us.es}
           {Dpto. de Ciencias de la Computación e Inteligencia Artificial} \\
      \href{http://www.us.es}
           {Universidad de Sevilla}  \\
    Sevilla, 10 de julio de 2023 (versión del \today)
  \end{tabular}\hfill\mbox{}}

%%%%%%%%%%%%%%%%%%%%%%%%%%%%%%%%%%%%%%%%%%%%%%%%%%%%%%%%%%%%%%%%%%%%%%%%%%%%%%
%% § Documento                                                              %%
%%%%%%%%%%%%%%%%%%%%%%%%%%%%%%%%%%%%%%%%%%%%%%%%%%%%%%%%%%%%%%%%%%%%%%%%%%%%%%

% \includeonly{}

% \includexmp{licencia}

\begin{document}

\maketitle
\newpage

% \input{Licencia/licenciaCC}
% \newpage

\tableofcontents
\clearpage

\section{Más sobre sucesiones}
\label{sec:org4af785d}

En la Lección 2, aprendimos la definición de una sucesión
\(a : \mathbb{N} \to \mathbb{R}\) que converge a un límite \(L\): para cualquier
tolerancia \(\epsilon > 0\), existe un momento \(N\), de modo que, para
cualquier punto posterior a ese momento, \(n \ge N\), estamos dentro de la
tolerancia, \(|a_n - L| < \epsilon\).

También aprendimos a hacer algo completamente trivial con ella:
demostrar que la sucesión constante converge, siendo esa constante su
límite. Vamos a subir un nivel, ¿os parece?

\section{Nivel 1: La propiedad arquimediana}
\label{sec:orgbf2f934}

La llamada propiedad arquimediana (que creo que se debe originalmente a
Eudoxo y aparece ya en el Libro V de los \emph{Elementos} de Euclides) es una
propiedad fundamental de los números reales que plasma la noción
intuitiva de que no existen números reales positivos «infinitamente
grandes» o «infinitamente pequeños».

Más concretamente, establece que no importa lo pequeño que sea \(\epsilon
> 0\),
siempre existe un número natural \(N\) tal que \(1/N\) es aún más
pequeño que \(\epsilon\) (y, por supuesto, positivo). De forma
equivalente, podemos enunciarlo como: para cualquier número real
positivo \(\epsilon\), existe un número natural \(N\) tal que \(1/\epsilon <
N\).

\textbf{¿Por qué es esto importante?} La propiedad arquimediana es una de las
propiedades más fundamentales que distinguen a los números reales de
otros sistemas numéricos. Sin ella, podríamos tener números positivos
«infinitamente grandes» o «infinitamente pequeños», lo que rompería la
mayor parte del cálculo y el análisis.

\begin{teorema}
Para cualquier \(\epsilon : \mathbb{R}\) con \(0 < \epsilon\), existe un \(N : \mathbb{N}\) tal que \(1/\epsilon < N\).
\end{teorema}

Esto es matemáticamente «obvio» para la mayoría de la gente: si tienes
un número positivo \(\epsilon\), por muy pequeño que sea, siempre puedes
encontrar un número natural lo suficientemente grande para que
\(1/\epsilon\) sea menor que él. Pero, ¿cómo se formaliza esto realmente
en Lean?

\subsection{Estrategia de demostración en lenguaje natural}
\label{sec:orge73f18e}

Primero, pensemos sobre esto en lenguaje natural. La idea clave es que
necesitamos proporcionar un número natural \(N\) específico que funcione.

Una opción natural sería utilizar algo relacionado con la función techo
(\emph{ceiling}). La función techo \(x \mapsto \lceil x \rceil\) redondea
cualquier número real al entero superior más cercano. Sin embargo, hay
un problema sutil: la función techo estándar toma valores en los enteros
\(\mathbb{Z}\), pero nosotros necesitamos valores en \(\mathbb{N}\) (los
números naturales). \textbf{¡¡No son la misma cosa!!}

Afortunadamente, Lean proporciona la «función techo para números
naturales» escrita como \(x \mapsto \lceil x \rceil_+\), que toma
cualquier número real y devuelve un número natural. (Puedes escribir
estos símbolos usando \texttt{\textbackslash{}lceil}, \texttt{\textbackslash{}rceil} y \texttt{\textbackslash{}\_+}). Para entradas negativas,
esta función devuelve 0. Por ejemplo, \(\lceil -3.14 \rceil_+ = 0\) y
\(\lceil 3.14 \rceil_+ = 4\).

Nuestra estrategia es:
\begin{itemize}
\item \textbf{Elección de \(N\):} Usar \(N = \lceil 1/\epsilon \rceil_+ + 1\)
\item \textbf{Por qué funciona:} Tenemos la «desigualdad clave»: \(1/\epsilon \le
  \lceil 1/\epsilon \rceil_+\), que se cumple por la definición de la
función techo.
\item \textbf{Obtención de la desigualdad estricta:} Al sumar 1 obtenemos \(1/\epsilon < \lceil 1/\epsilon \rceil_+ + 1\).
\end{itemize}

\subsection{Los desafíos de la implementación en Lean}
\label{sec:org7984489}

En Lean, nos encontramos con el problema de la \textbf{coerción de tipos} (\emph{type
coercion}, convertir entre diferentes tipos de números). \(\mathbb{N},
\mathbb{Z}, \mathbb{Q}\) y \(\mathbb{R}\) son todos tipos de cosas
diferentes.

Por ejemplo, al escribir:
\begin{verbatim}
have fact : 1 / ε ≤ ⌈1 / ε⌉₊ := by WhateverTheProofIs
\end{verbatim}

Lean lo registrará como:
\begin{verbatim}
fact : 1 / ε ≤ ↑⌈1 / ε⌉₊
\end{verbatim}

La flecha hacia arriba \texttt{↑} representa una función de coerción de los
naturales a los reales: \(↑ : \mathbb{N} \to \mathbb{R}\). Lean necesita
estas funciones explícitas para convertir entre tipos.

\subsection{Nuevas herramientas que necesitarás}
\label{sec:org5613c60}

\begin{itemize}
\item \texttt{⌈ · ⌉₊}: La función techo para números naturales.
\item \texttt{push\_cast}: Táctica que gestiona las coerciones entre tipos numéricos.
\item \texttt{bound}: Resuelve muchas desigualdades rutinarias y «triviales» una vez
alineados los tipos.
\end{itemize}

\subsection{Demostración en Lean}
\label{sec:org6430881}

\begin{verbatim}
lemma ArchProp {ε : ℝ} (hε : 0 < ε) :
∃ (N : ℕ), 1 / ε < N := by
  use ⌈ 1 / ε ⌉₊ + 1
  have fact : 1 / ε ≤ ⌈ 1 / ε ⌉₊ := by bound
  push_cast
  bound
\end{verbatim}

\subsection{Revisión de errores comunes}
\label{sec:orgfe1a633}

\begin{itemize}
\item ¡No uses la función techo regular \(\lceil \cdot \rceil\)! Devuelve
enteros (\texttt{Int}), no naturales (\texttt{Nat}).
\item \textbf{Cuidado con los problemas de casting:} si \texttt{bound} no funciona, intenta
usar primero \texttt{push\_cast}.
\end{itemize}
\end{document}
